\documentclass[11pt]{article}
\usepackage[margin=1in]{geometry}
\usepackage{graphicx}
\usepackage{amsmath, amssymb}
\usepackage{booktabs}
\usepackage{microtype}
\usepackage{caption}
\usepackage{xcolor}
\usepackage{hyperref}
\hypersetup{colorlinks=true, linkcolor=blue!50!black, urlcolor=blue!50!black}

\newcommand{\figfit}[1]{\includegraphics[width=\linewidth,height=0.75\textheight,keepaspectratio]{#1}}

\title{Book 3: Within-subject fine-grained analyses of the auditory
       superstitious-perception task}
\author{Auditory SP analyses}
\date{}

\begin{document}
\maketitle
\tableofcontents
\bigskip

%==============================================================================
\section{What this book is and how to read it}
%==============================================================================

Books 1 and 2 worked at the group level. Book 1 asked how the experimenter's
target / distractor labels relate to acoustic similarity between stimuli, and
Book 2 asked how subjects' responses relate to those labels and to one
another. The headline result of Book 2 was that group-level alignment is
small but not zero ($d'\approx 0.09$ in the \texttt{full\_sentence} block,
$p\!\approx\!0.006$), and that the cross-subject behavioural geometry has no
clear cluster structure.

Book 3 zooms inside individual subjects. The questions are:

\begin{enumerate}
\item \textbf{Per-subject classification image (Ch.~22).} If we average the
      raw waveforms of every stimulus a subject called ``target'' and
      subtract the average waveform of those they called ``distractor'',
      what does the resulting per-subject \emph{template} look like, and
      how aligned is it to the experimenter's true target-minus-distractor
      template (and to the original \texttt{wall.wav} sentence)?
\item \textbf{Within-subject internal consistency (Ch.~23).} Is each
      subject's behaviour reproducible across random halves of their
      trials, in either raw waveform or magnitude-spectrum space?
\item \textbf{Spectral logistic alignment (Ch.~24).} If we model each
      subject's responses as a linear logistic decision over 32 log-spaced
      magnitude-spectrum bands, how aligned is their weight vector to the
      analogous weight vector trained on the experimenter labels?
\item \textbf{Reminder-anchored and block-position effects (Ch.~25).} Every
      15 trials, subjects heard a \emph{reminder} replay of the
      \texttt{wall.wav} template. Do response, hit, and false-alarm rates
      depend on how many trials have elapsed since the last reminder,
      or on absolute position within the 150-trial block?
\item \textbf{Block-order and site (Ch.~26).} Does it matter whether a
      block was run first or second, and do Oberlin and UChicago
      participants differ?
\item \textbf{$d'$ tertiles -- who carries the signal (Ch.~27).} Split
      subjects within block by $d'$ and by $|d'|$ and ask which
      sub-group the alignment effect lives in.
\item \textbf{Individually-significant subjects (Ch.~28).} A close
      look at the few subjects who survive a per-subject permutation
      test, including their questionnaire profile.
\item \textbf{Questionnaire correlates (Ch.~29).} Do BAIS\_V (auditory
      imagery vividness), BAIS\_C (control), TAS, LSHS, DES, FSS,
      VHQ, or Stanford sleepiness predict any per-subject signal
      metric?
\end{enumerate}

\paragraph{Reading guide.}
Each chapter has four short sections: \emph{plain-English intent}, what we
do, the result you should look at, and how to interpret it. Every figure in
this book is a per-subject summary; group-level numbers are reported in the
text but the visual emphasis is always on within-subject distributions.

\paragraph{Reminder structure.}
Within each of the two 150-trial blocks the reminder screen fired before
trials $15, 30, 45, \ldots, 135$. So for any trial in either block, two
useful temporal covariates are
\(\texttt{trials\_since\_reminder} = \texttt{trial\_index} \bmod 15
   \in \{0,\ldots,14\}\)
and
\(\texttt{pos\_in\_block} = \texttt{trial\_index} \in \{0,\ldots,149\}\).

\paragraph{A useful hierarchical framing.}
Throughout this book each per-subject statistic $\theta_s$ (e.g.\ a cosine
alignment, a logistic slope) is implicitly viewed as
\(\theta_s \sim \mathcal{N}(\mu_\theta,\sigma_\theta^2)\) with
\(\mu_\theta\) the population effect and \(\sigma_\theta\) the
between-subject variability. A non-zero $\mu_\theta$ is what we mean by
``the typical subject is doing the task''; a positive $\sigma_\theta$ even
with $\mu_\theta=0$ would mean ``different subjects use systematically
different but real strategies''. We don't fit a full Bayesian hierarchical
model here, but every group-level statistic in this book is the
random-intercept version of the corresponding $\mu_\theta$.

%==============================================================================
\section{Chapter 22 -- Per-subject classification image}
%==============================================================================

\subsection{Plain-English intent}

For each subject, build the raw audio template they implicitly endorsed by
calling certain stimuli ``target'' and others ``distractor''. Compare it
to (a) the experimenter's own target-minus-distractor template for the
block, and (b) the original \texttt{wall.wav} clip the subject was told
to listen for.

\subsection{Method}

For subject $s$ in block $b$ with response vector
$y_s \in \{0,1\}^{150}$ (1 = called target):

\begin{align*}
\mathrm{CI}_s     &= \overline{w}_{\,y_{si}=1}
                   - \overline{w}_{\,y_{si}=0}, \\
\mathrm{CI}_{\rm true} &= \overline{w}_{\,t_{i}=1}
                          - \overline{w}_{\,t_{i}=0},
\end{align*}

where $w_i$ is the centred 2960-sample waveform of stimulus $i$ and
$t_i$ is its experimenter label. We report two cosines per subject:
$\cos(\mathrm{CI}_s,\,\mathrm{CI}_{\rm true})$ and
$\cos(\mathrm{CI}_s,\,\texttt{wall.wav})$. The null is built per subject
by permuting $y_s$ across stimuli ($N=1000$) and recomputing the cosine.

\subsection{Result}

Across both blocks the experimenter template $\mathrm{CI}_{\rm true}$ is
itself only mildly aligned to \texttt{wall.wav}
($\cos\!\approx\!0.42$): a target stimulus is acoustically a distorted,
embedded copy of the sentence rather than a clean copy.
At the per-subject level (Fig.~\ref{fig:ch22-hist}):

\begin{itemize}
\item \texttt{full\_sentence}: mean
      $\cos(\mathrm{CI}_s,\mathrm{CI}_{\rm true}) = +0.034$,
      \textbf{77\% of subjects above zero}, no individual subject
      permutation-significant at $p<0.05$.
\item \texttt{imagined\_sentence}: mean $+0.014$, 58\% above zero, one
      subject individually significant.
\end{itemize}

So the \emph{group} signal is real (a binomial test on $20/26$ positive
gives $p<0.01$ for \texttt{full\_sentence}) but no individual subject's
classification image is, on its own, reliably aligned. This is the
characteristic ``many small noisy effects, all in the same direction''
pattern. Figure~\ref{fig:ch22-examples} shows the worst- and
best-aligned subject in each block; even the best aligned $\mathrm{CI}_s$
does not look like a clean copy of $\mathrm{CI}_{\rm true}$.

\subsection{How to interpret}

Subjects appear to be \emph{slightly} biased towards endorsing
acoustically target-like stimuli, but the within-subject signal is
roughly the same order of magnitude as the per-subject permutation null.
What pushes the group result is the consistent \emph{sign} of the cosine
across subjects, not the magnitude in any one of them.

\begin{figure}[h]
\centering
\figfit{stimuli-analyses/outputs/book3/ch22-classification-image/ch22_alignment_hist.png}
\caption{Per-subject cosine between the subject's classification image
$\mathrm{CI}_s$ and the experimenter template $\mathrm{CI}_{\rm true}$.
Each bin is one subject. Vertical orange line: group mean.}
\label{fig:ch22-hist}
\end{figure}

\begin{figure}[h]
\centering
\figfit{stimuli-analyses/outputs/book3/ch22-classification-image/ch22_ci_examples.png}
\caption{Worst-aligned (top row) and best-aligned (bottom row) subject's
classification image (blue, normalised) overlaid on the experimenter
template (black, normalised), per block.}
\label{fig:ch22-examples}
\end{figure}

\begin{figure}[h]
\centering
\figfit{stimuli-analyses/outputs/book3/ch22-classification-image/ch22_ci_vs_wall_scatter.png}
\caption{Per-subject classification-image cosines: $y$-axis vs
$\mathrm{CI}_{\rm true}$, $x$-axis vs the original \texttt{wall.wav}.
Subjects above the horizontal axis are biased towards the experimenter
template; subjects to the right of the vertical axis are biased towards
the raw sentence.}
\label{fig:ch22-vs-wall}
\end{figure}

\clearpage

%==============================================================================
\section{Chapter 23 -- Within-subject split-half consistency}
%==============================================================================

\subsection{Plain-English intent}

If a subject is using \emph{any} consistent rule -- experimenter-aligned
or not -- to decide which stimuli to call ``target'', then the
classification image computed on a random half of their trials should
look like the classification image computed on the other half. We test
this in two feature spaces: raw centred waveform (2960-D) and the
magnitude FFT (1481-D).

\subsection{Method}

For each subject, repeat $200$ random splits of the 150 trials into
halves $A,B$. Compute $\mathrm{CI}_A,\mathrm{CI}_B$ on each half and
report the average Pearson correlation
$\bar r_s = \mathbb{E}_\text{splits}\;\mathrm{corr}(\mathrm{CI}_A,
                                                   \mathrm{CI}_B)$.
A response-shuffled null gives the floor.

\subsection{Result}

Both feature spaces give $\bar r_s$ values indistinguishable from zero
(group means in $[-0.007, +0.005]$) and indistinguishable from the
shuffled null (Fig.~\ref{fig:ch23-hist}). Plotting $\bar r_s$ against
the chapter-22 alignment cosine shows no relationship
(Fig.~\ref{fig:ch23-vsalign}).

\subsection{How to interpret}

Two readings:

\begin{enumerate}
\item Subjects are \emph{not} consistently using a rule that lives in raw
      waveform or magnitude-spectrum coordinates. Whatever consistency
      they have (chapter 22 says some sign-consistency exists) lives
      along the \emph{single} direction defined by
      $\mathrm{CI}_{\rm true}$, not along arbitrary directions in a
      $\sim$1500-D feature space. Estimating an arbitrary direction from
      75 trials is statistically hopeless even for perfectly consistent
      rules; estimating the projection onto a known direction is fine.
\item This is therefore a \emph{negative} result for ``hidden idiosyncratic
      strategies in raw acoustics''. If subjects each have their own
      private strategy in either feature space, the strategy is too weak
      to detect with 75 yes-trials and 75 no-trials per block.
\end{enumerate}

\begin{figure}[h]
\centering
\figfit{stimuli-analyses/outputs/book3/ch23-split-half/ch23_splithalf_hist.png}
\caption{Per-subject split-half consistency $\bar r_s$ in waveform
(top) and magnitude-spectrum (bottom) space, observed (blue) vs
response-shuffled (gray), per block. Both observed distributions are
indistinguishable from their shuffled controls.}
\label{fig:ch23-hist}
\end{figure}

\begin{figure}[h]
\centering
\figfit{stimuli-analyses/outputs/book3/ch23-split-half/ch23_consistency_vs_alignment.png}
\caption{No relationship between within-subject consistency
($y$, ch.~23) and label-template alignment ($x$, ch.~22).}
\label{fig:ch23-vsalign}
\end{figure}

\clearpage

%==============================================================================
\section{Chapter 24 -- Per-subject spectral logistic regression}
%==============================================================================

\subsection{Plain-English intent}

Reduce each stimulus to its magnitude-spectrum amplitude in 32 log-spaced
frequency bands ($\sim$50\,Hz to $\sim$3950\,Hz), then fit a ridge
logistic regression per subject:
\(P(\text{called\_target}_i = 1) = \sigma(w_s^\top x_i + b_s)\).
Compare the resulting weight vector $w_s$ to the analogous weight
vector $w_{\rm true}$ trained on the experimenter labels for the
same block.

This is the same spirit as chapter 22, but in a much lower-dimensional
basis where 75 yes-trials \emph{is} enough to estimate a direction.

\subsection{Method}

Per stimulus, compute the magnitude FFT and average within each of 32
geometrically-spaced bands. Standardise across the 150 stimuli.
Per-subject ridge logistic with $C=0.5$ (sklearn,
\texttt{liblinear}). $w_{\rm true}$ is the same model fit with
$y$=experimenter label. Per-subject permutation null: shuffle $y_s$ and
refit ($N=200$ shuffles).

\subsection{Result}

\begin{itemize}
\item \texttt{full\_sentence}: mean
      $\cos(w_s,w_{\rm true}) = +0.064$, 65\% positive,
      \textbf{2/26 subjects individually $p<0.05$}.
\item \texttt{imagined\_sentence}: mean $+0.042$, 54\% positive, no
      subject individually significant.
\end{itemize}

Figures \ref{fig:ch24-hist} and \ref{fig:ch24-heatmap-fs} show the
distribution and the per-subject weight vectors sorted by alignment.
In \texttt{full\_sentence} the highest-aligned subjects' weight
vectors visibly mirror the dominant low-frequency structure of
$w_{\rm true}$.
Figure~\ref{fig:ch24-vs-dprime} shows that within \texttt{full\_sentence}
subjects with higher chapter-16 $d'$ also tend to have larger
$\cos(w_s,w_{\rm true})$ -- exactly what one would expect if
the spectral-band logistic is capturing the same per-subject
discrimination signal that SDT captures.

\subsection{How to interpret}

This is the strongest within-subject evidence in the book that
individual participants are using the experimenter-defined acoustic
information, not just at the group level. The effect is small in
magnitude but consistent in direction, and the spectral basis is small
enough that two individual subjects survive a per-subject permutation
test.

\begin{figure}[h]
\centering
\figfit{stimuli-analyses/outputs/book3/ch24-spectral-logistic/ch24_alignment_hist.png}
\caption{Per-subject cosine alignment between the subject's spectral
logistic weight vector $w_s$ and the experimenter weight vector
$w_{\rm true}$. Dashed gray line: subject-averaged null mean.}
\label{fig:ch24-hist}
\end{figure}

\begin{figure}[h]
\centering
\figfit{stimuli-analyses/outputs/book3/ch24-spectral-logistic/ch24_weight_heatmap_full_sentence.png}
\caption{\texttt{full\_sentence}: $w_{\rm true}$ across the 32 bands
(top) and per-subject $w_s$ rows (bottom), sorted by alignment cosine.
$x$-axis tick labels: band centre frequency in Hz.}
\label{fig:ch24-heatmap-fs}
\end{figure}

\begin{figure}[h]
\centering
\figfit{stimuli-analyses/outputs/book3/ch24-spectral-logistic/ch24_weight_heatmap_imagined_sentence.png}
\caption{\texttt{imagined\_sentence}: same as previous figure for the
imagined-sentence block.}
\label{fig:ch24-heatmap-is}
\end{figure}

\begin{figure}[h]
\centering
\figfit{stimuli-analyses/outputs/book3/ch24-spectral-logistic/ch24_alignment_vs_dprime.png}
\caption{Per-subject spectral alignment vs SDT sensitivity $d'$
(chapter 16). Subjects more sensitive in the SDT sense also tend
to have spectral weights that align with the experimenter's.}
\label{fig:ch24-vs-dprime}
\end{figure}

\clearpage

%==============================================================================
\section{Chapter 25 -- Reminder-anchored and block-position effects}
%==============================================================================

\subsection{Plain-English intent}

Every 15 trials a reminder screen replayed the wall.wav template. Three
related questions:
(A) Do response, hit, and false-alarm rates differ between the first 5
trials after a reminder, the middle 5, and the last 5 (just before the
next reminder)?
(B) After regressing out the experimenter label, is there a per-subject
slope of $P(\text{yes})$ on either trials-since-reminder or absolute
position in the block, and is the average of those slopes non-zero?
(C) Does the block-level mean response rate drift across the 10
reminder segments?

\subsection{Method}

Per-subject logistic regression of $\texttt{called\_target}$ on
$z(\texttt{trials\_since\_reminder})$,
$\texttt{true\_target}$,
$z(\texttt{trials\_since\_reminder})\times\texttt{true\_target}$,
$z(\texttt{pos\_in\_block})$, and
$z(\texttt{pos\_in\_block})\times\texttt{true\_target}$,
then group-level Wilcoxon of each per-subject slope against zero.

\subsection{Result}

\begin{itemize}
\item No reliable bin-of-15 effect: in both blocks the mean
      response/hit/FA rates over first-5 / middle-5 / last-5 bins
      essentially overlap (Fig.~\ref{fig:ch25-bin5}).
\item No reliable trials-since-reminder slope (\texttt{full\_sentence}
      Wilcoxon $p=0.69$; \texttt{imagined\_sentence} $p=0.76$) and no
      reliable absolute-position slope ($p=0.39$ and $p=0.92$).
\item The interaction $\texttt{trials\_since\_reminder}
      \times\texttt{true\_target}$ is also null in both blocks.
\item The within-subject coefficient on $\texttt{true\_target}$ itself
      \emph{is} reliably positive in \texttt{full\_sentence}
      (median $+0.17$, Wilcoxon $p=0.011$) and in the same direction
      but not significant in \texttt{imagined\_sentence}
      (median $+0.04$, $p=0.53$). This is the per-subject SDT signal,
      now seen as a non-zero population effect after partialling out
      temporal covariates.
\item Block-level segment-drift slopes are centred near zero in both
      blocks (Fig.~\ref{fig:ch25-drift}). No general fatigue or warm-up
      pattern in the response rate.
\end{itemize}

\subsection{How to interpret}

The reminder cadence does not modulate behaviour in any detectable way.
There is no warm-up after the reminder, no decay before the next one,
and no drift across the block. The only temporal covariate that matters
in the per-subject regression is the experimenter's label itself: the
typical subject's $P(\text{yes})$ is significantly higher on true
targets than on true distractors, replicating the chapter-16
group-level finding from a different angle (mixed-effects-style
per-subject estimation rather than collapsed counts).

\begin{figure}[h]
\centering
\figfit{stimuli-analyses/outputs/book3/ch25-reminder-anchored/ch25_bin5_group.png}
\caption{Group-level mean response, hit, false-alarm, and accuracy
rates as a function of the 5-trial bin since the most recent reminder.
Error bars: SEM across subjects.}
\label{fig:ch25-bin5}
\end{figure}

\begin{figure}[h]
\centering
\figfit{stimuli-analyses/outputs/book3/ch25-reminder-anchored/ch25_slope_forest.png}
\caption{Per-subject logistic-regression coefficients for the five
temporal/label covariates. Gray points: individual subjects; blue
point: group mean $\pm$ SE. The only coefficient with a reliably
non-zero population mean is \texttt{beta\_truetarget} in
\texttt{full\_sentence}.}
\label{fig:ch25-forest}
\end{figure}

\begin{figure}[h]
\centering
\figfit{stimuli-analyses/outputs/book3/ch25-reminder-anchored/ch25_segment_drift.png}
\caption{Per-subject slope of mean response rate across the 10
reminder segments (10$\times$15-trial bins). No reliable drift.}
\label{fig:ch25-drift}
\end{figure}

\clearpage

%==============================================================================
\section{Synthesis}
%==============================================================================

What Book 3 found, end-to-end:

\begin{enumerate}
\item In the \texttt{full\_sentence} block, the typical subject's
      classification image is weakly but consistently aligned with the
      experimenter's true target-minus-distractor template
      ($20/26$ positive cosines, $p<0.01$ binomial). No individual
      subject is alignment-significant on a per-subject permutation,
      but the population mean is non-zero.
\item In a 32-band log-spectral basis, the same subjects'
      ridge-logistic weight vectors align with the experimenter's
      ($+0.064$ mean cosine, $65\%$ positive, two individually
      significant subjects). Per-subject spectral alignment is
      positively correlated with chapter-16 $d'$.
\item Within-subject consistency in arbitrary high-dimensional
      waveform or spectrum coordinates is at chance: subjects do
      \emph{not} appear to use idiosyncratic acoustic strategies that
      are detectable from 75 yes-trials per block.
\item Reminder-anchored time-course is flat. No first-5 / middle-5 /
      last-5 differences, no trials-since-reminder slope, no drift
      across the 10 reminder segments. The reminder cadence is
      behaviourally invisible.
\item The only reliable per-subject temporal-regression effect is the
      coefficient on the true experimenter label
      (\texttt{full\_sentence} median $+0.17$,
      Wilcoxon $p=0.011$): the typical subject is more likely to say
      ``yes'' on true-target trials. This is a within-subject
      restatement of the chapter-16 group $d'$.
\end{enumerate}

\paragraph{Where the signal lives.}
The picture that emerges across Books 2 and 3 is consistent: there is a
small, positive, group-level alignment with the experimenter labels,
which can be isolated either via SDT counts (Ch.~16), classification
images projected onto the true template (Ch.~22), spectral logistic
weights (Ch.~24), or a per-subject random-intercept logistic with the
true label as a predictor (Ch.~25). All four routes give signed,
non-zero population effects in \texttt{full\_sentence}; all four are
much weaker (and rarely significant) in \texttt{imagined\_sentence}.

\paragraph{Where the signal does not live.}
Idiosyncratic per-subject strategies in raw acoustic feature spaces are
not detectable. Reminder-anchored time-course effects are not
detectable. Block-position drift is not detectable. Within-subject
internal consistency in raw acoustic spaces is at chance.

\paragraph{Implicit hierarchical model.}
Every per-subject statistic in this book has been treated as
\(\theta_s \sim \mathcal{N}(\mu_\theta,\sigma_\theta^2)\). The data
are consistent with $\mu_\theta>0$ for the alignment statistics
in \texttt{full\_sentence} and with $\mu_\theta\approx 0,
\sigma_\theta\approx 0$ (no strategy variation beyond noise) for
all the temporal covariates. A full Bayesian hierarchical model
\(y_{si} \sim \mathrm{Bernoulli}(\sigma(w_s^\top x_i + b_s))\) with
$w_s\sim\mathcal{N}(\mu_w,\Sigma_w)$ would be the natural next step
if a sharper estimate of $\mu_\theta$ relative to $\sigma_\theta$
were needed; here the random-intercept Wilcoxon and per-subject
permutation tests already make the qualitative point.

\clearpage

%==============================================================================
\section{Chapter 26 -- Block-order and site effects}
%==============================================================================

\subsection{Plain-English intent}

Each subject ran two blocks (\texttt{full\_sentence} and
\texttt{imagined\_sentence}) in a counterbalanced order. Two questions:
(A) does it matter whether a block was run first or second
(practice / fatigue / familiarisation with \texttt{wall.wav})? and
(B) do the two recruitment sites (Oberlin, $N=16$; UChicago, $N=10$)
differ on any of our headline metrics?

\subsection{Method}

For each per-subject metric (Ch.16 $d'$, Ch.22 $\cos(\mathrm{CI}_s,
\mathrm{CI}_{\rm true})$, Ch.24 $\cos(w_s,w_{\rm true})$, Ch.25
\texttt{beta\_truetarget}) we (i) compare Block-1 vs Block-2 within
\texttt{block\_type} via Mann--Whitney, and (ii) compare the same
metric across all 26 subjects via a paired Wilcoxon on
$(\theta_{\text{run-second}}-\theta_{\text{run-first}})$ irrespective
of which scheme came second. We also run Mann--Whitney by site within
each block.

\subsection{Result}

\paragraph{Block order.}
Across all four metrics the Block-2 mean is consistently larger than
the Block-1 mean (Fig.~\ref{fig:ch26-blockorder}). The paired
Wilcoxon trend is strongest for \texttt{beta\_truetarget}
($\bar{\Delta}=+0.16$, $W=104$, $p=0.071$) and
$d'$ ($\bar{\Delta}=+0.09$, $p=0.099$). No comparison reaches
$p<0.05$ but the consistent direction of the four trends suggests a
mild practice/familiarisation effect.

\paragraph{Site.}
Numerically Oberlin tracks higher than UChicago on every metric in
every block (e.g.\ \texttt{full\_sentence} $d'$: Oberlin $0.135$ vs
UChicago $0.021$; cos\_ci\_true: $0.057$ vs $-0.002$). None of the
eight Mann--Whitney tests reach $p<0.05$ ($p$-values in $[0.12,0.69]$)
so this remains descriptive. Figure~\ref{fig:ch26-site} shows the
distributions side-by-side.

\begin{figure}[h]
\centering
\figfit{stimuli-analyses/outputs/book3/ch26-block-order-and-site/ch26_block_order.png}
\caption{Per-subject metrics by block-order index, within block type.
Black bars: group means. Block-2 means are uniformly higher than
Block-1 means.}
\label{fig:ch26-blockorder}
\end{figure}

\begin{figure}[h]
\centering
\figfit{stimuli-analyses/outputs/book3/ch26-block-order-and-site/ch26_site_split.png}
\caption{Per-subject metrics by site (Oberlin vs UChicago). Oberlin
tracks higher on every metric in every block but no comparison is
$p<0.05$.}
\label{fig:ch26-site}
\end{figure}

\clearpage

%==============================================================================
\section{Chapter 27 -- $d'$ tertiles: who carries the signal}
%==============================================================================

\subsection{Plain-English intent}

The group means in chapters 22 and 24 are small. If they are driven by
a few subjects with genuinely good performance, splitting subjects into
$d'$ tertiles should reveal it. We split (within each block) into
\emph{signed} $d'$ tertiles (low $\le$ med $\le$ high) and \emph{absolute}
$|d'|$ tertiles (worst $\le$ middle $\le$ most-extreme). Kruskal--Wallis
tests whether each per-subject statistic differs across tertiles.

\subsection{Result -- signed $d'$}

Almost every metric is monotonically ordered low--med--high in both
blocks (Kruskal $p<10^{-4}$ for $d'$ itself, $p<10^{-3}$ for
\texttt{cos\_ci\_true} in \texttt{imagined\_sentence}, etc.). This is
unsurprising and largely a consistency check.

\subsection{Result -- absolute $|d'|$}

The interesting split. In \texttt{full\_sentence} every alignment
metric is significantly elevated in the highest $|d'|$ tertile:

\begin{center}
\begin{tabular}{lrrrr}
\toprule
metric & low $|d'|$ & med $|d'|$ & high $|d'|$ & Kruskal $p$ \\
\midrule
$d'$              & $+0.015$ & $+0.042$ & $+0.210$ & $0.005$ \\
\texttt{cos\_ci\_true}     & $+0.011$ & $+0.009$ & $+0.081$ & $0.026$ \\
\texttt{cos\_w\_true}      & $-0.001$ & $+0.022$ & $+0.165$ & $0.230$ \\
\texttt{beta\_truetarget}  & $+0.038$ & $+0.081$ & $+0.341$ & $0.006$ \\
\texttt{accuracy}          & $0.503$  & $0.508$  & $0.541$  & $0.006$ \\
\bottomrule
\end{tabular}
\end{center}

The high-$|d'|$ tertile in \texttt{full\_sentence} is also high
on \emph{signed} $d'$ (mean $+0.21$, all positive), so high-$|d'|$
subjects are not high-$|d'|$ because of \emph{anti-aligned}
performance -- they are simply the few subjects who actually
discriminated the experimenter targets. In
\texttt{imagined\_sentence}, by contrast, $|d'|$ tertile is
\emph{not} predictive of any alignment metric ($p>0.32$ for
\texttt{cos\_ci\_true}, \texttt{cos\_w\_true},
\texttt{beta\_truetarget}). Imagined-block subjects with extreme
$|d'|$ are split between aligned and anti-aligned -- consistent with
the chapter-16 group $d'\!\approx\!0$.

\begin{figure}[h]
\centering
\figfit{stimuli-analyses/outputs/book3/ch27-dprime-tertiles/ch27_signed_tertile_box.png}
\caption{Per-subject metrics by signed $d'$ tertile. Monotonic across
all metrics in both blocks (largely a consistency check).}
\label{fig:ch27-signed}
\end{figure}

\begin{figure}[h]
\centering
\figfit{stimuli-analyses/outputs/book3/ch27-dprime-tertiles/ch27_abs_tertile_box.png}
\caption{Per-subject metrics by absolute $|d'|$ tertile. In
\texttt{full\_sentence} the high-$|d'|$ tertile carries the
alignment signal; in \texttt{imagined\_sentence} the highest-$|d'|$
subjects are mixed in sign and do not differ in alignment from the
other tertiles.}
\label{fig:ch27-abs}
\end{figure}

\clearpage

%==============================================================================
\section{Chapter 28 -- Individually-significant subjects}
%==============================================================================

\subsection{Plain-English intent}

A per-subject permutation test with $N=1000$--$200$ shuffles can
detect subjects whose alignment is unlikely to arise from chance.
Three tests:
$d'$ from chapter 16,
$\cos(\mathrm{CI}_s,\mathrm{CI}_{\rm true})$ from chapter 22, and
$\cos(w_s,w_{\rm true})$ from chapter 24.

\subsection{Result}

8 subject-blocks (out of $26\times 2=52$) survive a one-sided
$p<0.05$ on at least one test (3 in \texttt{full\_sentence},
5 in \texttt{imagined\_sentence}). No subject is significant in
\emph{both} blocks. The per-test counts:

\begin{center}
\begin{tabular}{lccc}
\toprule
block & sig on $d'$ & sig on CI alignment & sig on spectral alignment \\
\midrule
\texttt{full\_sentence}     & 0 & 2 & 3 \\
\texttt{imagined\_sentence} & 1 & 3 & 3 \\
\bottomrule
\end{tabular}
\end{center}

A weaker but more pervasive marker is being \emph{stably aligned} on
all three signal metrics within a block (cos\_ci\_true $>0$ AND
cos\_w\_true $>0$ AND beta\_truetarget $>0$): 15 / 26 in
\texttt{full\_sentence} and 10 / 26 in \texttt{imagined\_sentence}.

\paragraph{Questionnaire profile of the significant subjects.}
The most striking pattern is in BAIS\_V (Bucknell Auditory Imagery
Scale -- vividness). 7 of the 8 significant subjects scored at or
above mid-range on BAIS\_V; 6 of the 8 scored $\ge 68$ (close to
ceiling, range $14$--$98$). The single significant subject with floor
BAIS\_V ($=14$, ``no auditory imagery'') is significant only on the
spectral test, not on the imagery-relevant CI test.

\begin{figure}[h]
\centering
\figfit{stimuli-analyses/outputs/book3/ch28-significant-subjects/ch28_signif_overlap_venn.png}
\caption{Per-test counts of individually-significant subjects in each
block.}
\label{fig:ch28-counts}
\end{figure}

\begin{figure}[h]
\centering
\figfit{stimuli-analyses/outputs/book3/ch28-significant-subjects/ch28_significant_metric_table.png}
\caption{Per-significant-subject metric values, normalised within
column. Rows labelled \texttt{site\_initial subject\_id\_tail
block\_type\_prefix}.}
\label{fig:ch28-table}
\end{figure}

\clearpage

%==============================================================================
\section{Chapter 29 -- Questionnaire correlates}
%==============================================================================

\subsection{Plain-English intent}

For each subject we have eight questionnaire scores: BAIS\_V (auditory
imagery vividness), BAIS\_C (auditory imagery control), VHQ (visual
imagery), TAS (Tellegen absorption), LSHS (Launay--Slade hallucination
proneness), DES (dissociative experiences), FSS (flow state), and
SSS\_mean (Stanford Sleepiness, mean across pre/post measurements).
We compute Spearman $\rho$ with each performance metric, per block.

\subsection{Method}

Spearman $\rho$ across all $8 \times 8 \times 2 = 128$
quest$\times$perf$\times$block combinations, with simple Bonferroni
correction. We also run a one-sided BAIS\_V median-split test
(prediction: high BAIS\_V $\Rightarrow$ higher alignment).

\subsection{Result}

No correlation survives Bonferroni at $p<0.05$ (which would require
$p<3.9\times 10^{-4}$). Three uncorrected $p<0.05$ correlations:

\begin{center}
\begin{tabular}{llcrr}
\toprule
block & quest$\times$perf & $n$ & $\rho$ & $p$ \\
\midrule
\texttt{imagined\_sentence} & DES vs criterion        & 26 & $+0.44$ & $0.023$ \\
\texttt{full\_sentence}     & BAIS\_C vs hit\_rate    & 26 & $+0.40$ & $0.042$ \\
\texttt{imagined\_sentence} & DES vs FA rate          & 26 & $-0.40$ & $0.045$ \\
\bottomrule
\end{tabular}
\end{center}

\paragraph{BAIS\_V median split.}
The pre-registered direction (high BAIS\_V $\Rightarrow$ higher
within-subject alignment) is consistent across all four
\texttt{full\_sentence} metrics, with one-sided $p$-values
$\in [0.056, 0.075]$:

\begin{center}
\begin{tabular}{lrrr}
\toprule
metric & high BAIS\_V (N=13) & low BAIS\_V (N=13) & one-sided $p$ \\
\midrule
$d'$                  & $0.152$ & $0.030$ & $0.056$ \\
\texttt{cos\_ci\_true}     & $0.061$ & $0.008$ & $0.069$ \\
\texttt{cos\_w\_true}      & $0.085$ & $0.042$ & $0.341$ \\
\texttt{beta\_truetarget}  & $0.240$ & $0.072$ & $0.076$ \\
\bottomrule
\end{tabular}
\end{center}

\paragraph{How to interpret.}
Each individual test is marginal, but four of four \texttt{full\_sentence}
metrics point in the predicted direction with $p<0.10$, and the
chapter-28 categorical observation (7 of 8 individually-significant
subject-blocks have above-mid BAIS\_V) is consistent. This is the
strongest cross-questionnaire signal in the dataset.

\begin{figure}[h]
\centering
\figfit{stimuli-analyses/outputs/book3/ch29-questionnaires/ch29_bais_v_scatter.png}
\caption{BAIS\_V vs each within-subject signal metric, coloured by
site. Spearman $\rho$ and uncorrected $p$ in each panel title.
Dashed line: OLS fit.}
\label{fig:ch29-baisv}
\end{figure}

\begin{figure}[h]
\centering
\figfit{stimuli-analyses/outputs/book3/ch29-questionnaires/ch29_bais_c_scatter.png}
\caption{BAIS\_C (auditory imagery control) vs each within-subject
signal metric.}
\label{fig:ch29-baisc}
\end{figure}

\begin{figure}[h]
\centering
\figfit{stimuli-analyses/outputs/book3/ch29-questionnaires/ch29_tas_scatter.png}
\caption{TAS (Tellegen absorption) vs each within-subject signal
metric.}
\label{fig:ch29-tas}
\end{figure}

\begin{figure}[h]
\centering
\figfit{stimuli-analyses/outputs/book3/ch29-questionnaires/ch29_lshs_scatter.png}
\caption{LSHS (Launay--Slade hallucination proneness) vs each
within-subject signal metric.}
\label{fig:ch29-lshs}
\end{figure}

\clearpage

%==============================================================================
\section{Synthesis (revised)}
%==============================================================================

Putting chapters 22--29 together:

\begin{enumerate}
\item \textbf{Group-level alignment is small but real in
      \texttt{full\_sentence}} (Ch.~22, $20/26$ positive cosines;
      Ch.~24, $+0.064$ mean spectral alignment; Ch.~25, true-target
      coefficient median $+0.17$, $p=0.011$).
\item \textbf{The signal is concentrated in the high-$|d'|$ tertile}
      (Ch.~27): in \texttt{full\_sentence}, the top third of
      $|d'|$-ranked subjects has cos\_ci\_true $\approx +0.08$ and
      beta\_truetarget $\approx +0.34$, vs.\ $\approx 0$ in the lower
      two thirds. The same split is null in
      \texttt{imagined\_sentence}.
\item \textbf{8 of 52 subject-blocks are individually permutation-
      significant} (Ch.~28). No subject is significant in both
      blocks, but $25/52$ are stably aligned (positive on all three
      signal metrics) within a block.
\item \textbf{Block-2 means trend higher than Block-1 means} on every
      metric (Ch.~26), suggesting practice / template-familiarisation,
      with paired-Wilcoxon $p\approx 0.07$ for the strongest case.
\item \textbf{Site differences are descriptive only} (Ch.~26):
      Oberlin $>$ UChicago numerically on every metric, no test
      $p<0.05$.
\item \textbf{Reminder cadence is behaviourally invisible} (Ch.~25):
      no first-5 / mid-5 / last-5 effect, no
      trials-since-reminder slope, no segment drift.
\item \textbf{High auditory imagery vividness (BAIS\_V) tracks
      alignment} (Ch.~28, Ch.~29): 7 of 8 individually-significant
      subjects are above mid BAIS\_V, and a one-sided median-split
      test in \texttt{full\_sentence} gives $p<0.08$ on three of
      four signal metrics. No questionnaire correlation survives
      Bonferroni, so this is suggestive rather than confirmed.
\item \textbf{Idiosyncratic high-D acoustic strategies are not
      detectable} (Ch.~23). Subjects who carry the alignment signal
      do so along the experimenter-defined direction, not along
      arbitrary directions.
\end{enumerate}

\paragraph{One-sentence summary.}
Within the \texttt{full\_sentence} block a minority of subjects --
preferentially those with high auditory-imagery vividness -- align
their responses to the experimenter's acoustic targets in a way that
is detectable through four mutually consistent statistics
(SDT $d'$, classification image alignment, spectral logistic
alignment, true-target coefficient). The \texttt{imagined\_sentence}
block does not show this pattern, the reminder cadence is irrelevant,
and within-subject idiosyncratic strategies are below detection.

\end{document}

